\documentclass[11pt]{article}\usepackage{layout}\begin{document}
\begin{canvas}\Header{第二章 · 全球开源生态地图}{teal}\Footer{01}
\Box{15}{27}{180}{\fontsize{21}{27}\selectfont\color{ink}\sffamily \bfseries 2.8 全球开源行业地图}
\Box{15}{40}{180}{\fontsize{9.7}{13.5}\selectfont\color{muted}\sffamily 开源服务于哪些行业？本节借鉴北美行业分类体系（NAICS），按文本估计仓库应用领域。2026 年 11 月 9 日快照覆盖 18,086 个英文文本样本，呈现行业结构的延续与 AI 信号的分化。}
\Box{15}{53}{180}{\includegraphics[width=180mm]{.build/assets/fig-2-8-1.pdf}}
\Box{15}{199}{180}{\fontsize{7.8}{10.8}\selectfont\color{muted}\sffamily 图 2.8-1  19 个行业的位次演变与 2025 创建组规模。来源：团队 OpenDigger–NAICS 研究及配套数据。年份为仓库创建年份，文本来自同一期后快照；每个仓库归入一类。}
\Box{15}{211}{87}{\fontsize{10.5}{14.2}\selectfont\color{teal}\sffamily \bfseries ① 宏观格局延续，中部位次交替}
\Box{15}{217}{87}{\fontsize{10.1}{15.6}\selectfont\color{ink}2021 与 2025 两端的行业位次相关系数为 0.982，信息业始终居首；中部行业则有交替。这表明，AI 相关表述的增加，可以发生在总体相似的行业结构之中。}
\Box{108}{211}{87}{\fontsize{10.5}{14.2}\selectfont\color{teal}\sffamily \bfseries ② 开源供给已跨越单一信息行业}
\Box{108}{217}{87}{\fontsize{10.1}{15.6}\selectfont\color{ink}2025 创建组中，信息业占 17.5\%，其余 82.5\% 分布于教育、专业服务、金融、制造及公共服务等领域。行业标签由此补充了技术分类，让开源的潜在服务对象更清晰。}
\Box{15}{245}{87}{\fontsize{10.5}{14.2}\selectfont\color{teal}\sffamily \bfseries ③ 教育与文娱的相对位置上移}
\Box{15}{251}{87}{\fontsize{10.1}{15.6}\selectfont\color{ink}教育由第 3 位升至第 2 位，份额由 10.9\% 变为 12.3\%；艺术娱乐休闲由第 7 位升至第 6 位，份额由 6.5\% 变为 7.8\%。它们提示开源供给具有广泛的知识与内容应用基础。}
\Box{108}{245}{87}{\fontsize{10.5}{14.2}\selectfont\color{teal}\sffamily \bfseries ④ 位次须与份额结合判断}
\Box{108}{251}{87}{\fontsize{10.1}{15.6}\selectfont\color{ink}金融与保险由第 5 位升至第 4 位，份额却从约 8.2\% 变为 8.1\%；公用事业在第 3—7 位之间波动。相对位置变化不能直接解释为行业规模的增减，更不能替代实际使用评价。}
\end{canvas}
\begin{canvas}\Header{第二章 · 全球开源生态地图}{teal}\Footer{02}
\Box{15}{27}{180}{\fontsize{21}{27}\selectfont\color{ink}\sffamily \bfseries 2.8 AI 提及：总体分化与行业落差}
\Box{15}{40}{180}{\fontsize{9.7}{13.5}\selectfont\color{muted}\sffamily 行业格局较为稳定，AI 相关表述却在不同创建群体和行业间呈现明显差异。这里的“AI 提及”指简介、标签和说明文本出现相关术语，是技术表述信号，不能直接等同于实际采用率。}
\Box{15}{53}{180}{\includegraphics[width=180mm]{.build/assets/fig-2-8-2.pdf}}
\Box{15}{199}{180}{\fontsize{7.8}{10.8}\selectfont\color{muted}\sffamily 图 2.8-2  AI 提及的创建群体差异、词表对照及行业比较。来源：团队配套统计表；区间为 Wilson 95\% 区间，仅反映抽样波动。右图按比例排序，行业数量与分类可靠性需一并考虑。}
\Box{15}{211}{87}{\fontsize{10.5}{14.2}\selectfont\color{teal}\sffamily \bfseries ① 较新仓库中的 AI 表述更普遍}
\Box{15}{217}{87}{\fontsize{10.1}{15.6}\selectfont\color{ink}宽词表下，2021 创建组为 6.1\%，2025 创建组为 24.2\%，相差 18.1 个百分点，约为前者的四倍。这是同一快照中的群体差异，反映技术表述变化，尚不能定位实际采用时间。}
\Box{108}{211}{87}{\fontsize{10.5}{14.2}\selectfont\color{teal}\sffamily \bfseries ② 行业落差大于总体平均所示}
\Box{108}{217}{87}{\fontsize{10.1}{15.6}\selectfont\color{ink}金融与保险为 40.1\%，公用事业为 8.1\%，相差约 32 个百分点。专业科技服务达到 47.3\%，但分类误差较大；建筑、行政支持等小样本行业的高比例也应谨慎解读。}
\Box{15}{245}{87}{\fontsize{10.5}{14.2}\selectfont\color{teal}\sffamily \bfseries ③ 收窄词表，主要信号依然存在}
\Box{15}{251}{87}{\fontsize{10.1}{15.6}\selectfont\color{ink}窄词表去除泛化机器学习用语后，2025 创建组仍达到 22.8\%，与宽词表相差约 1.4 个百分点。较新创建组的高提及比例，不完全由宽泛术语造成，但仍需实际功能与部署证据。}
\Box{108}{245}{87}{\fontsize{10.5}{14.2}\selectfont\color{teal}\sffamily \bfseries ④ AI 信号并不局限于信息行业}
\Box{108}{251}{87}{\fontsize{10.1}{15.6}\selectfont\color{ink}制造、医疗与社会援助分别为 29.1\%、28.9\%，均高于全样本 24.2\%；信息业为 22.7\%。这为后续行业案例提供线索：应结合任务、数据与运行条件，核查技术表述如何转化为应用。}
\end{canvas}
\begin{canvas}\Header{第二章 · 全球开源生态地图}{teal}\Footer{03}
\Box{15}{27}{180}{\fontsize{21}{27}\selectfont\color{ink}\sffamily \bfseries 2.8 行业结构与 AI 集中度的错位}
\Box{15}{40}{180}{\fontsize{9.7}{13.5}\selectfont\color{muted}\sffamily 行业规模、AI 提及比例与 AI 相关仓库数量，是三个不同的问题。将它们放在一起，才能看见哪些领域集中了更多技术表述，以及哪些行业需要进一步核查应用与复用条件。}
\Box{15}{53}{180}{\includegraphics[width=180mm]{.build/assets/fig-2-8-3.pdf}}
\Box{15}{199}{180}{\fontsize{7.8}{10.8}\selectfont\color{muted}\sffamily 图 2.8-3  AI 相关仓库的行业构成、份额差及创建群体差异。来源：团队汇总表与论文图 3，重绘方法见附录 1。面积与份额以 836 个 AI 提及仓库为分母；＊专业科技分类待复核。}
\Box{15}{211}{87}{\fontsize{10.5}{14.2}\selectfont\color{teal}\sffamily \bfseries ① 三个行业集中了过半 AI 仓库}
\Box{15}{217}{87}{\fontsize{10.1}{15.6}\selectfont\color{ink}专业科技、信息、金融合计 435 个 AI 提及仓库，占 52.0\%；三者占全部仓库的份额为 36.9\%。AI 表述具有比仓库分布更强的行业集中性，值得进一步追踪其共享工具与知识资源。}
\Box{108}{211}{87}{\fontsize{10.5}{14.2}\selectfont\color{teal}\sffamily \bfseries ② 规模优势与 AI 集中度并不一致}
\Box{108}{217}{87}{\fontsize{10.1}{15.6}\selectfont\color{ink}信息业仓库最多，但 AI 提及数量为 137 个，低于专业科技服务的 186 个。专业科技、金融的 AI 仓库份额分别比自身仓库份额高 10.9、5.3 个百分点；观察需求应结合规模和结构。}
\Box{15}{245}{87}{\fontsize{10.5}{14.2}\selectfont\color{teal}\sffamily \bfseries ③ 行业群体差异并非单向递增}
\Box{15}{251}{87}{\fontsize{10.1}{15.6}\selectfont\color{ink}制造、医疗在 2023 创建组的比例约为 23.7\%、23.5\%，高于 2024 组的 15.4\%、14.7\%，2025 组又较高。群体构成和文本更新均可能产生这种差异，不能直接视为应用曾经退潮。}
\Box{108}{245}{87}{\fontsize{10.5}{14.2}\selectfont\color{teal}\sffamily \bfseries ④ 把差异转化为能力建设的问题}
\Box{108}{251}{87}{\fontsize{10.1}{15.6}\selectfont\color{ink}两端提及率差距从公用事业的 3.4 个百分点，到专业科技的 38.5 个百分点不等。公共投入宜据行业任务，核查数据、成本、本地化与维护条件；以可复用能力和实际结果检验普惠价值。}
\end{canvas}
\begin{canvas}\Header{附录 1 · 行业地图口径与来源}{teal}\Footer{A1}
\Box{15}{28}{180}{\fontsize{20}{27}\selectfont\color{ink}\sffamily \bfseries 附录 1｜行业地图口径与证据}
\Box{15}{46}{180}{\fontsize{11.3}{16}\selectfont\color{teal}\sffamily \bfseries 1  观察对象、时间与分类含义}
\Box{15}{55}{180}{\fontsize{9.4}{14.5}\selectfont\color{ink}本节为期后专题，不并入 2025.07—2026.06 主统计期。2026-11-09 从 OpenDigger 所跟踪的公开、非 Fork（非平台复制）仓库中，按创建年份各抽取 4,000 个；候选池合计 269,767 个。以文本前 2,000 字符的 ASCII 比例大于 90\% 近似筛选英文，保留 18,086 个。未逐仓核验开源许可。NAICS 原用于企业经营单位，本研究将其迁移为仓库应用领域标签；覆盖 19 类，不含企业管理类（55）。}
\Box{15}{90}{180}{\fontsize{11.3}{16}\selectfont\color{teal}\sffamily \bfseries 2  分类可靠性与图表复核}
\Box{15}{99}{180}{\fontsize{9.4}{14.5}\selectfont\color{ink}采用 NAICS-GH 的 5,270 个训练及验证样本训练词语与字符特征的线性分类器，独立测试 1,318 个。准确率 76.3\%，加权 F1 为 75.8\%；F1 同时衡量误分与漏分。专业科技服务 F1 为 0.54，金融为 0.90，公用事业为 0.88，各类可靠性不一，且测试集表现不等于本研究样本准确率。图 2.8-3 的 30 个分年值由附件矢量曲线恢复，并与整数分母及 2025 汇总表核对；原始行业分年明细仍待补归档。}
\Box{15}{134}{180}{\fontsize{11.3}{16}\selectfont\color{teal}\sffamily \bfseries 3  AI 提及的含义与统计边界}
\Box{15}{143}{180}{\fontsize{9.4}{14.5}\selectfont\color{ink}在简介、标签与 README（项目说明）的合并文本前 1,800 字符中识别 AI／机器学习关键词；窄词表去掉泛化机器学习用语。文本来自同一快照，年份只代表创建群体，不能定位实际采用时间、推断生成式 AI 的因果影响或外推到 2030。行业占比的 Wilson 95\% 区间不覆盖分类错误、词表误差及样本偏差。}
\Box{15}{173}{180}{\fontsize{11.3}{16}\selectfont\color{teal}\sffamily \bfseries 4  证据来源}
\Box{15}{182}{180}{\fontsize{8.8}{12.5}\selectfont\color{muted}\sffamily [S1] 团队研究稿：Stable Structure, Uneven Diffusion: A Reproducible Industry-Level Study of AI Mentions in Open Source；配套《开源仓库行业版图数据报告》及复现包，2026-11-09 快照。正文数值优先采用复现包汇总表。}
\Box{15}{203}{180}{\sffamily\fontsize{8.8}{12.5}\selectfont\color{muted}[S2] Xu, K. \& Quispe, A. (2026). \href{https://arxiv.org/abs/2607.06505}{Industry Classification of GitHub Repositories Using the North American Industry Classification System (NAICS). arXiv:2607.06505.}}
\Box{15}{218}{180}{\sffamily\fontsize{8.8}{12.5}\selectfont\color{muted}[S3] \href{https://huggingface.co/datasets/aquiro1994/naics-gh}{NAICS-GH 数据集（6,588 个带标签仓库）}；[S4] \href{https://www.census.gov/naics/}{美国人口普查局：NAICS 官方说明}。外部资料核查于 2026-11-10。}
\Box{15}{236}{180}{\fontsize{11.3}{16}\selectfont\color{teal}\sffamily \bfseries 5  版本、待核查项与评审重点}
\Box{15}{245}{180}{\fontsize{9.4}{14.5}\selectfont\color{ink}v0.2 重绘行业位次与规模、AI 总体与行业差异、AI 仓库构成与热力矩阵三幅综合图；压缩方法性配图与说明卡片。待核查：行业分年原始明细、专业科技及小样本分类、多语覆盖和实际复用案例。请重点评审宏观结构与行业差异是否更清楚，图表信息取舍是否恰当；暂不将探索性行业排序提升为全面领先或采用率结论。}
\end{canvas}
\end{document}